Neutrino quasielastic scattering $\nu_{l}N(n) \rightarrow l^{-}p$ on nuclei~$N$ 
is a dominant signal process for neutrino oscillation experiments
and is used to extract information about the axial vector form factor for 
nucleons~\cite{Gran:2006jn,AguilarArevalo:2007ab,Espinal:2007zz,Lyubushkin:2008pe,Nakajima:2011zz}.
The simple final-state topology combined with an assumption that the initial-state 
nucleon is at rest allows for an estimate of the neutrino energy 
from the final-state lepton kinematics alone. This estimate can be altered 
by the fact that the nucleon is bound in a nucleus, which is often modeled 
by assuming that the nucleons are non-interacting within a relativistic 
Fermi gas (RFG). However, measurements based on the final-state lepton on nuclei 
with $A$~$>$~2 over different ranges of four-momentum transfer \QSq are inconsistent 
with the RFG model with $M$$_{A}$~$\sim$~1~GeV in the absolute size of the cross section per nucleon 
and the predicted energy deposition near the vertex of these 
interactions~\cite{Gran:2006jn,AguilarArevalo:2007ab,Nakajima:2011zz,AguilarArevalo:2013hm,Fields:2013zhk,Fiorentini:2013ezn}. 

Many models of nuclear effects attempt to explain these discrepancies by considering  
possible correlations between nucleons. These include short-range correlations 
as observed in electron scattering~\cite{Benhar:2006wy,Arrington:2011xs,Arrington:2012ax}, 
long-range correlations that are modeled with the random phase approximation 
(RPA)~\cite{Graczyk:2003ru,Nieves:2004wx,Valverde:2006yi,Martini:2009uj,Martini:2010ex}, 
and meson exchange currents 
(MEC)~\cite{Martini:2009uj,Martini:2010ex,Nieves:2011pp,Bodek:2011ps,Shen:2012xz,Sobczyk:2012ms,Gran:2013kda}.
Each of these processes changes the event rate and final-state particle
kinematics.

In addition, hadrons produced in neutrino-nucleus interactions can undergo final-state 
interactions (FSI) as they propagate through the nucleus. Consequently, 
a sample including only a lepton and nucleons will invariably  
contain events from inelastic processes. These include $\Delta(1232)$ resonance 
production and decay, where the pion is not observed. Events from both 
inelastic processes with no final-state pions and nucleon-nucleon correlations 
contribute to the measured quasielastic (QE) cross section, but have different kinematics and 
final-state hadron content for the same neutrino energy. Therefore, these events 
can alter the accuracy of any neutrino energy 
estimate~\cite{Leitner:2010kp,Martini:2012fa,Martini:2012uc,Nieves:2012yz,Lalakulich:2012gm,Lalakulich:2012hs,Mosel:2013fxa} 
that neutrino oscillation experiments~\cite{LBNE,Jediny:2014lda} use.  

Additional information is accessible through measurements of the 
hadronic component. Previous 
measurements~\cite{Gran:2006jn,Espinal:2007zz,Lyubushkin:2008pe,Nakajima:2011zz} 
have made a selection on the energy and/or direction of a tracked proton and 
its consistency with the QE hypothesis in order to increase the QE purity of 
the sample. Such a selection will remove events modified by FSI or caused by 
non-QE processes. The presented QE-like analysis specifically retains 
sensitivity to these effects by requiring only that the final-state proton's 
direction and momentum be measured.

Presented is a differential cross-section measurement of QE-like events
that consist of a muon with at least one proton and no pions in the final state. 
By using the kinetic energy of the most energetic (leading) proton, a measurement 
of \QSq is made from the hadronic component alone. This extracted cross section 
is measured with much improved acceptance at large muon scattering angles and 
higher \QSq than that of \minerva's QE measurements that rely on muon 
kinematics~\cite{Fields:2013zhk,Fiorentini:2013ezn}. 

